\newpage
\section{Purification, sonship and life together}
\zlabel{triptych:concise:commentary:start}
\subsection{How can those receiving holiness purify themselves?}
First John places the Father's love before the believer's action. The children receive their identity and then purify themselves because they hope to see him. The difficulty concerns agency: the same person receives holiness and does something about becoming holy. Augustine's explanation begins by asking whether the name ``child'' describes a real life. A physician's title would be useless without the ability to heal. Likewise sonship must be real in faith, hope and charity. The promised future concerns people whose adoption has truly begun.\footnote{1 John 3:1--3; Augustine, \work{Tractates on First John} IV.4--5.}

Augustine answers the apparent self-purification directly. God purifies, and the person consents to his action. God does not purify an unwilling subject. Asking for help shows both the need and the action of the person asking: a helper assists someone who also acts. Consent consequently belongs within the gift's reception. The Collect's request for mercy can be understood in this way. The assembly's petition is an act of people who acknowledge their need and turn toward God; the saints' intercession accompanies that petition.\footnote{Augustine, \work{First John} IV.7.}

Aquinas encounters the same question from the other direction in the psalm. Innocence appears before blessing and mercy. A reader might understand the order as a reward given after someone has made himself innocent. Aquinas permits that account of reward, but also offers a causal explanation: God's blessing and mercy make the person's innocence and access possible. The order of sentences can describe the life of the recipient without making him the independent cause of what he receives. Aquinas's argument concerns the psalm's promise; Augustine's concerns the will named in John's command.\footnote{Aquinas, \work{Super Psalmos} 23, n. 4, paragraph 86991.}

Bede's explanation of Revelation identifies the cause of cleansing in another image. The white robes have been washed in the Lamb's blood. He distinguishes the martyrs' blood from the blood of Jesus that cleanses the whole Church. This answers a different danger: identifying the saving cause with what the saints have suffered. His comment permits the multitude to include saints who were not killed as martyrs, while retaining faithful confession and tribulation. The received Victorinus commentary also understands the garments through baptism and the preservation of received grace.\footnote{Bede, \work{Explanation of the Apocalypse}, I, on 7:14, p. 53; received Victorinus, \work{Commentary on the Apocalypse}, seventh chapter, comment 9. The latter is cited as the received text, with its recension and attribution difficulties.}

Augustine and Aquinas explain how the person becomes capable of the promised communion. Bede explains the cleansing shared by the people who stand before God. In the moral sense, the difference directs attention both to the will's cooperation and to the neighbor who shares the need for mercy.

\subsection{What changes in a person who hopes to see God?}
Augustine explains preparation with two images that address different needs. A container intended for a large gift must be stretched. A vessel filled with vinegar must be emptied and cleaned before receiving honey. The first comparison explains a growing capacity; the second explains the removal of incompatible desires. Greater intensity alone would not purify desire if the person continued to seek happiness in the same lesser attachments. Nor do honey, wine or gold describe God's substance: Augustine says the comparisons fall short of what is promised.\footnote{Augustine, \work{First John} IV.6.}

This helps explain the psalm's reference to vanity. Augustine describes a soul seeking lasting eternity instead of treating itself as a passing possession. Innocent hands keep that desire connected to deeds. In Aquinas's explanation, exterior innocence and interior purity are distinguishable demands. His warning about pride in one's apparent purity adds a further difficulty: a person may become attached to his own moral condition. The desire for God requires the heart to relinquish that self-satisfaction as well as more obvious disordered appetites.\footnote{Augustine, \work{Expositions on the Psalms}, Psalm 23.4; Aquinas, \work{Super Psalmos} 23, n. 3, paragraph 86990.}

Aquinas follows the promise of vision into a question about its object. Scripture calls God unseen, yet Jesus promises the pure heart that it will see God. A created reflection of divine brightness would leave John's promise of seeing him as he is unfulfilled. Aquinas also reasons from knowledge and desire: someone who knows effects seeks their cause. A final happiness that left the first cause unknown would leave that desire unfinished. He therefore understands the promise as sight of God himself.\footnote{Aquinas, \work{Super Matthaeum}, cap. 5, lectio 2, paragraph 87340, on Matthew 5:8, adducing 1 John 3:2 and 1 Corinthians 13:12.}

The mode of that sight is a separate question. Aquinas rejects bodily vision of the divine essence because God is not a colored object. Jesus promises sight to the heart; Aquinas explains it intellectually. Augustine makes a related distinction through Christ's humanity and divinity. Those condemned at judgment can see Christ's human form. The sight promised with likeness and purification must therefore exceed that encounter: Augustine speaks of the divine Word, equal with the Father. The promised likeness remains creaturely likeness, even in its perfection.\footnote{Augustine, \work{First John} IV.5,9; Aquinas, \work{Super Matthaeum}, paragraph 87340, on 5:8.}

Here the anagogical sense, concerning final fulfillment, is more specific than a hope for improvement. The person hopes to see God and love him in completed communion. The moral sense concerns acts preparing for that end. Aquinas distinguishes those acts from consummated happiness while allowing present beginnings of blessedness. The Prayer after Communion's petition for charity's completion can be understood beside this distinction. The sacrament has truly been received; the recipient still seeks an end beyond his present condition.

The communal interpretation asks what that end means for relations with others. The Prayer over the Offerings speaks of the saints' continuing concern for those seeking salvation. Fulfillment in God remains compatible with such care. Augustine's vessel images explain a recipient's preparation, whereas this prayer's subject concerns the perfected saints' solicitude. Their combination allows desire for vision to be understood within a communion where persons continue to love those in need. The allegorical sense identifies that communion with Christ's people, cleansed through him; the promised fulfillment includes those persons together.

\subsection{How are the numbered and innumerable people related?}
The seer hears a number and then sees a multitude. The appointed passage preserves that transition while omitting the twelve tribal counts between the scenes. The number is exact, the crowd beyond human enumeration. Bede explains the number through twelve multiplied by twelve and a thousand. Twelve refers to the apostles and to faith in the Trinity throughout the fourfold world. The Church's fullness is his conclusion. Patriarchal descent can concern natural lineage or imitation of faith. The explanation therefore extends membership beyond ancestry alone.\footnote{Revelation 7:2--4,9--14; Bede, \work{Explanation}, I, on 7:4, pp. 45--46.}

His explanation of verse 9 admits two ways of following the narrative. It may resume the description of victorious glory after recapitulation. It may also proceed from Israel, where the Gospel was first preached, to the salvation of the Gentiles. In the latter account the succession of groups has a historical rationale in the order of preaching. Bede thus understands the number as fullness while leaving open whether this transition distinguishes Israel and the Gentiles. Neither explanation makes the first number a maximum population of heaven.\footnote{Bede, \work{Explanation}, I, on 7:9, p. 52.}

The multitude's common act is worship of God and the Lamb. Its members still have nations and languages in the description. Their unity concerns their shared relation to the Savior. Augustine's explanation of the psalm gives a compatible account of the nations believing in the glorified Lord as his Church. Aquinas calls the seekers of God one generation. Taken together, these accounts explain a people whose common origin in grace calls for a common manner of life: innocent hands, a pure heart and desire for God's face.\footnote{Augustine, \work{Expositions}, Psalm 23.1--6; Aquinas, \work{Super Psalmos} 23, n. 4.}

The people before the throne have emerged from tribulation. Bede connects palms with the triumph of the Cross and victory in Christ. He retains entry through many tribulations and gives particular attention to the final persecution, without dating it. The saints' victory remains Christ's gift and includes their fidelity. In the interpretation centered on purified desire, that endurance tests the will's dependence on grace. In the interpretation centered on the gathered people, it describes the condition through which a community preserves its confession.

Bede and the received Victorinus commentary agree about cleansing but differ about the eastern angel. The received Victorinus text identifies Elias; Bede identifies the incarnate Lord, the messenger of the Father's will.\footnote{Bede, \work{Explanation}, I, on 7:2,9,14, pp. 44--45,52--53; received Victorinus, \work{Commentary}, seventh chapter, comments 2,9.}

\subsection{What kind of peace resembles the Son?}
Augustine describes peace through an ordering within the person: desires are subject to reason and reason to God. It is a condition of the heart that has ceased to let disordered appetite rule. Chrysostom explains peacemaking through the activity of the Son, who unites the divided and reconciles the alienated. A disciple doing that work bears the Son's likeness. One account identifies the order needed within the person; the other identifies work done between persons.\footnote{Augustine, \work{Our Lord's Sermon on the Mount} I.2.9; Chrysostom, \work{Homilies on Matthew} XV.7.}

Chrysostom makes the social demand exact. Avoiding one's own quarrels leaves something undone when other people remain estranged. The peacemaker must help reconcile them. Aquinas includes interior order, peace with others and this further work among those at odds. A calm person may still neglect that work. Conversely, someone attempting reconciliation needs his own desires ordered if pride and resentment are to cease directing the effort. These are complementary explanations of charity with different objects, rather than competing accounts in which one kind of peace cancels the other.\footnote{Aquinas, \work{Super Matthaeum}, paragraph 87340, on Matthew 5:9.}

John's present sonship gives a further reason for conduct that resembles the Son. The Father's love makes children before their work displays their likeness. In the broader chapter, love is expressed in deeds and help to a brother in need. The activity follows a received identity. This is compatible with Chrysostom's explanation of peacemakers: the promised name identifies a likeness displayed in reconciliation. In the communal interpretation, Bede explains the people Christ cleanses, while Chrysostom explains the reconciling activity they learn from the Son.\footnote{1 John 3:1--3,16--18; verses 16--18 are contextual, beyond the appointment.}

Jesus follows peacemaking with persecution for justice. Chrysostom explains the adjacency by denying that every peace deserves blessing. An agreement may abandon virtue or the assistance another person needs. His account of justice includes godliness, virtue and help to others. The Gospel consequently gives grounds for refusing an agreement that would require those goods to be surrendered. It also preserves the command to reconcile; the possibility of resistance does not remove the duty to seek another's good.\footnote{Chrysostom, \work{Matthew} XV.7.}

The next verses put limits on the interpretation of hostility. Christ names false accusations and suffering for his sake. Chrysostom and Aquinas both require attention to the cause. Injury alone cannot establish the injured person's righteousness. Chrysostom also explains the special wound of slander: danger can sometimes bring public praise, but false reports remove even that consolation. His account gives endurance a concrete object while retaining the tests of truth and fidelity to Christ.\footnote{Chrysostom, \work{Matthew} XV.7--9; Aquinas, \work{Super Matthaeum}, paragraphs 87564--87566.}

Augustine relates the world's failure to recognize God's children to Christ's own unrecognized lowliness and reproof of sin. Public approval therefore cannot determine whether adoption is real. Chrysostom's conditions prevent that explanation from becoming permission to disregard every criticism. In the moral sense, the believer must examine both the charge and the conduct that produced it. At Communion, purity, peacemaking and persecution are named in one selection. Hope for the kingdom includes these distinct demands on the communicant's conduct.\footnote{Augustine, \work{First John} IV.4; Matthew 5:8--10.}

\subsection{Which promises are already being fulfilled?}
The Beatitudes speak both of a kingdom already possessed and of gifts still promised. Augustine explains their order as conversion through humility, teachable meekness, mourning, hunger for justice, mercy, purity and peace. His text places meekness before mourning, unlike the U.S. proclaimed sequence. His account of mercy gives dependence a practical consequence: someone asking God's help may be stronger than a neighbor in a particular matter and should give the help he can. Need therefore becomes a reason for merciful action.\footnote{Augustine, \work{Sermon on the Mount} I.3--4, §§10--12.}

The promised earth receives materially different explanations. Augustine understands a stable eternal inheritance. Chrysostom includes a present sensible reward alongside future gifts. He argues that violent striving can cost a person property and life, while meekness can preserve them. Aquinas reports Chrysostom and prefers a future account. Chrysostom's argument concerns the consequences of meekness; it supplies no guarantee that every meek person will retain his possessions. Augustine's future interpretation likewise allows blessedness to have real beginnings now.\footnote{Augustine, \work{Sermon on the Mount} I.2.4; Chrysostom, \work{Matthew} XV.5; Aquinas, \work{Super Matthaeum}, paragraph 87340.}

The invitation in the acclamation helps explain how costly conduct can belong to rest. Chrysostom considers the yoke in Matthew 11:29--30 beside Christ's demanding commands. Meekness and humility change how the disciple bears it by relieving the burdens of pride and sin. Those called are already burdened by anxiety, sorrow or sin. The call offers help within an exacting life. Beside Augustine's course, this explanation gives humility a place both in beginning to learn and in continuing to act.\footnote{Matthew 11:25--30; Chrysostom, \work{Matthew} XXXVIII.3.}

The Mass's prayers concern both assistance now and a shared destination. The Preface describes Jerusalem as mother and speaks of the saints' help and example for pilgrims. Its Jerusalem language belongs to a composed prayer, without adding a biblical appointment. The Prayer after Communion asks for perfected charity and the heavenly banquet. These subjects allow the two interpretations to reach different, compatible conclusions. Preparation for vision concerns the person's capacity to know and love God; reconciliation concerns persons arriving together in a charity that benefits others.

In the allegorical sense, both conclusions concern life in the Son: the Lamb's cleansing makes the Church his people, and peacemakers resemble his reconciling work. In the anagogical sense, present communion awaits completion in the sight of God and the common banquet. The saints' continuing intercession is compatible with that fulfillment. The assembly still petitions for grace after receiving Communion. Its hope includes perseverance and growth rather than an assurance that present reception has already completed its journey.
